\chapter{PySpark Deep Dive and Performance Tuning}
\index{PySpark}\index{performance tuning}\index{AQE}\index{broadcast join}\index{sort-merge join}%
\index{salting}\index{pandas UDF}\index{shuffle partitions}\index{Ganglia}%
\begin{objectives}
\begin{itemize}
    \item Explain how Catalyst and Tungsten compile and execute DataFrame code, including whole-stage codegen.
    \item Choose join strategies---broadcast hash versus sort-merge---and direct them with hints and AQE.
    \item Tune shuffle partitions, memory, and parallelism for a given data volume.
    \item Remediate skew with salting and AQE skew handling, with measured evidence.
    \item Replace Python UDFs with built-ins or pandas UDFs and quantify the difference.
    \item Deliver an evidence-based performance report with attributable, isolated improvements.
\end{itemize}
\end{objectives}

\begin{crosscloud}
Spark tuning is \textbf{portable}---the same knobs exist on every managed Spark; what changes is
the default and the telemetry surface.

\vspace{2mm}
{\footnotesize
\begin{tabular}{@{}p{3.0cm}p{3.4cm}p{3.2cm}p{3.2cm}@{}}
\toprule
\textbf{Capability} & \textbf{Databricks} & \textbf{AWS (EMR/Glue)} & \textbf{Azure / GCP} \\
\midrule
Adaptive execution & AQE on by default & AQE (EMR 6.x+) & AQE (Synapse/Fabric/Dataproc) \\
Execution engine & Photon (vectorized) & EMR runtime & Fabric native / vanilla Spark \\
Observability & Spark UI + query profile + Ganglia & Spark UI + CloudWatch & Spark UI + Synapse/Fabric monitoring \\
Join hints & \texttt{broadcast()}, SQL hints & Same API & Same API \\
Python UDF accel & pandas UDFs, Arrow & pandas UDFs & pandas UDFs \\
Skew handling & AQE skew split + salting & AQE skew (EMR) & AQE skew \\
\bottomrule
\end{tabular}}

\vspace{1mm}
Outside Spark: \textbf{Snowflake} pushes the same work into its engine via Snowpark with warehouse
sizing as the main lever; \textbf{BigQuery} abstracts execution entirely---you trade the knobs for
simplicity and lose the tuning surface this chapter teaches.
\end{crosscloud}

\section{Week 9 Roadmap and Mental Model}
Chapter 4 taught the execution model and the UI; Week 9 turns that literacy into a tuning method. The promise: most Spark performance problems are one of five---a bad join strategy, wrong shuffle partition count, skew, an expensive UDF, or a missing cache/predicate---and each has a signature you can read in the plan or the UI \cite{sparkSQLTuning,chambersZaharia}.

The mental model: \textbf{tune in a loop with isolated variables}. Measure a baseline, change exactly one thing, measure again, attribute. A report that changes four configs at once proves nothing and teaches nobody; the acceptance criteria of this week's lab enforce the discipline.

\begin{definitionbox}
\textbf{Adaptive Query Execution (AQE)} re-optimizes a query at stage boundaries using runtime statistics: coalescing small shuffle partitions, switching sort-merge joins to broadcast, and splitting skewed partitions. It is on by default in the Databricks Runtime; tuning means knowing when its defaults misfit.
\end{definitionbox}

\begin{keyterms}
\textbf{Whole-stage codegen}: fusing operators into one generated Java/native function. \textbf{Broadcast join}: replicating the small side to every executor, avoiding the shuffle. \textbf{Salting}: artificially spreading a hot key across buckets. \textbf{pandas UDF}: a vectorized UDF exchanging Arrow columnar batches instead of per-row Python calls. \textbf{Spill}: intermediate data evicted to disk under memory pressure.
\end{keyterms}

\section{Monday: Catalyst, Tungsten, and When Not to RDD}
\subsection{The Optimization Stack}
Your DataFrame code becomes an unresolved logical plan, then an analyzed plan (catalog-resolved), then an optimized logical plan (predicate pushdown, column pruning, constant folding, join reordering by statistics), then a physical plan (join strategy, aggregation mode), then generated code fused into whole stages. Reading \texttt{explain(True)} output top-to-bottom is reading this pipeline.

\subsection{When (Not) to Use RDDs}
RDDs trade the optimizer for control: no Catalyst, no codegen, no Photon, manual partitioning. The valid reasons in 2026 are rare---custom partitioner algorithms, legacy libraries. Every workload in this program stays on the DataFrame/SQL API; the moment you write \texttt{.rdd}, you have left the performance warranty.

\begin{pitfall}
\textbf{Collecting to the driver ``to inspect.''} \texttt{collect()}, \texttt{toPandas()}, and Python loops over rows move data to the driver; at scale they are the out-of-memory classic. Inspect with \texttt{show}/\texttt{limit}/aggregates; if you need local analysis, sample explicitly with \texttt{.sample()} first.
\end{pitfall}

\section{Tuesday: Joins and AQE}
\subsection{Broadcast versus Sort-Merge}
A \textbf{sort-merge join} shuffles both sides by the join key, sorts, and merges---correct at any size, expensive by construction. A \textbf{broadcast hash join} sends the small side to every executor and probes a hash table---no shuffle of the big side. The threshold (\texttt{spark.sql.autoBroadcastJoinThreshold}, default 10 MB) and AQE's runtime statistics decide; you direct with \texttt{broadcast(dim\_df)} when you know better than the estimates \cite{sparkSQLTuning}.

\begin{lstlisting}[language=Python,caption={Directing the join strategy with evidence.},label={lst:ch9-joins}]
from pyspark.sql import functions as F

facts = spark.read.table("de_training.gold.fact_orders")
prods = spark.read.table("de_training.gold.dim_product")

# If dim_product is small (check!), force broadcast:
facts.join(F.broadcast(prods), "product_id").explain(True)
# Physical plan should show BroadcastHashJoin, no Exchange on the fact side.

# AQE will convert SMJ to BHJ at runtime when stats prove it fits;
# the query profile shows "AdaptiveSparkPlan" with the final strategy.
\end{lstlisting}

\subsection{Skew and Salting}
Skew defeats parallelism: one task holds 80\% of a shuffle. AQE's skew-splitting helps by splitting the hot partition after the fact; \textbf{salting} prevents it by exploding the hot key across N buckets and replicating the small side per bucket. Chapter 4's lab built the muscle; this week's discipline is knowing when each applies.

\section{Wednesday: Shuffle Tuning, UDFs, and Memory}
\subsection{Shuffle Partitions and I/O}
\texttt{spark.sql.shuffle.partitions} (default 200) sets post-shuffle parallelism. Target $\sim$100--200 MB per partition: 200 GB shuffled $\rightarrow$ $\sim$1500 partitions. With AQE coalescing, overshooting is safer than undershooting. Watch \emph{spill} in the UI---spilling tasks signal memory pressure that partition sizing or executor memory must absorb.

\subsection{UDF Performance}
A row-at-a-time Python UDF serializes every row between the JVM and Python---often 10--100x slower than a native expression. The ladder, best to worst: built-in SQL functions, \textbf{pandas (Arrow) UDFs} over columnar batches, plain Python UDFs, RDD \texttt{map}.

\begin{lstlisting}[language=Python,caption={The UDF ladder: native, then vectorized, then---rarely---plain Python.},label={lst:ch9-udfs}]
from pyspark.sql import functions as F
from pyspark.sql.functions import pandas_udf
import pandas as pd

# Best: built-in
df.withColumn("amount_band", F.when(F.col("amount") > 100, "high")
                              .otherwise("low"))

# Good: vectorized pandas UDF when the logic truly needs Python
@pandas_udf("string")
def risky_band(s: pd.Series) -> pd.Series:
    return pd.cut(s, bins=[0, 50, 200, 1e9], labels=["low", "mid", "high"]).astype(str)

df.withColumn("risk_band", risky_band("amount"))

# Last resort: plain Python UDF (row-at-a-time serialization cost)
\end{lstlisting}

\begin{tipbox}
Before writing any UDF, search the function list: \texttt{spark.sql.functions} covers dates, strings, arrays, maps, math, and windowing. Most ``I need a UDF'' moments are undiscovered built-ins.
\end{tipbox}

\section{Thursday Hands-on Project: The Attributable Tuning Report}
Take a deliberately degraded 2 TB-scale job (or a 200 GB surrogate) and produce a tuning report where every improvement is isolated and measured.
\index{hands-on lab}\index{salting!lab}\index{broadcast join!lab}

\begin{lstlisting}[language=Python,caption={Week 9 lab: four isolated optimizations, one evidence table.},label={lst:ch9-lab}]
from pyspark.sql import functions as F
import time

def timed(label, fn):
    t0 = time.time(); fn(); return (label, round(time.time() - t0, 1))

facts = spark.read.table("de_training.silver.big_facts")
dim   = spark.read.table("de_training.silver.big_dim")

results = []
# Baseline: shuffle join, default partitions, Python UDF
spark.conf.set("spark.sql.adaptive.enabled", "false")
results.append(timed("baseline", lambda:
    facts.join(dim, "k").groupBy("dim.v").agg(F.sum("amount")).count()))

# Opt 1: right-size shuffle partitions for the measured volume
spark.conf.set("spark.sql.shuffle.partitions", "1500")
results.append(timed("shuffle_partitions_1500", lambda:
    facts.join(dim, "k").groupBy("dim.v").agg(F.sum("amount")).count()))

# Opt 2: broadcast the measured-small dimension
results.append(timed("broadcast_dim", lambda:
    facts.join(F.broadcast(dim), "k")
         .groupBy("dim.v").agg(F.sum("amount")).count()))

# Opt 3: re-enable AQE (skew split + coalesce) and compare
spark.conf.set("spark.sql.adaptive.enabled", "true")
results.append(timed("aqe_on", lambda:
    facts.join(F.broadcast(dim), "k")
         .groupBy("dim.v").agg(F.sum("amount")).count()))

spark.createDataFrame(results, ["experiment", "seconds"]) \
     .write.mode("overwrite").saveAsTable("de_training.silver.week09_results")
\end{lstlisting}

\subsection*{Deliverables}
\begin{itemize}
    \item The evidence table: one row per experiment, identical result values across all.
    \item Spark UI screenshots per experiment highlighting the stage that changed.
    \item A written report: symptom, hypothesis, change, measurement, conclusion---per optimization.
\end{itemize}

\subsection*{Acceptance Criteria}
\begin{itemize}
    \item Each experiment changes exactly one variable; confounded results are redone.
    \item Result correctness is reconciled across all experiments.
    \item The report names the single largest win and explains its mechanism from the plan.
    \item One experiment deliberately regresses (e.g., broadcasting the large side) and is documented as a negative control.
\end{itemize}

\subsection*{Stretch Goals}
\begin{itemize}
    \item Add a salted-join variant for an injected hot key and compare against AQE skew handling.
    \item Replace a Python UDF with a pandas UDF and quantify rows/second processed.
    \item Repeat the best configuration on Photon off vs.\ on (runtime toggle) and quantify the engine effect.
\end{itemize}

\section{Friday Use Case: The 2 TB Join That Misses Its SLA}
A logistics partner's daily join job (2 TB facts against 40 GB of reference data) must finish by 06:00; it currently finishes at 09:15. You are engaged for a one-week performance engagement with a written deliverable.
\index{SLA}\index{performance engagement}

\subsection{Requirements and Assumptions}
The job runs on a fixed-size job cluster; the cluster may be resized only with justification. Reference data grows 5\% monthly. The business accepts logic-preserving rewrites but no result changes.

\subsection{Architecture Decisions}
\textbf{The engagement structure.} Day 1: capture three consecutive runs' Spark UI exports and \texttt{explain} plans---the baseline dossier. Day 2: attribute time per stage; expect the classic finding (the ``small'' reference side grew past the broadcast threshold months ago, flipping the join to sort-merge). Days 3--4: isolated fixes with A/B runs---forced broadcast with a memory check, shuffle partition sizing, skew audit on the tenant key. Day 5: regression suite (row-for-row reconciliation against the slow baseline), the report, and the handover runbook.

\begin{figure}[H]
    \centering
    \resizebox{\textwidth}{!}{%
    \begin{tikzpicture}[
        ph/.style={rectangle, rounded corners, draw=primary, thick, fill=primary!7, align=center, minimum width=2.9cm, minimum height=1.0cm, font=\small\bfseries},
        art/.style={rectangle, rounded corners, draw=codegreen!60!black, thick, fill=green!7, align=center, minimum width=2.9cm, minimum height=0.8cm, font=\scriptsize},
        flow/.style={-{Latex[length=2mm]}, draw=primary, thick}
    ]
        \node[ph] (base) at (0,0) {Baseline\\dossier};
        \node[ph] (attr) at (3.6,0) {Attribute\\per stage};
        \node[ph] (fix)  at (7.2,0) {Isolated\\fixes};
        \node[ph] (reg)  at (10.8,0) {Regression\\reconcile};
        \node[ph] (rep)  at (14.4,0) {Report +\\runbook};
        \node[art] (a1) at (0,-1.5) {UI exports, plans};
        \node[art] (a2) at (3.6,-1.5) {stage time table};
        \node[art] (a3) at (7.2,-1.5) {A/B measurements};
        \node[art] (a4) at (10.8,-1.5) {row-for-row diff};
        \node[art] (a5) at (14.4,-1.5) {SLA projection};
        \draw[flow] (base) -- (attr);
        \draw[flow] (attr) -- (fix);
        \draw[flow] (fix) -- (reg);
        \draw[flow] (reg) -- (rep);
        \foreach \a/\b in {base/a1, attr/a2, fix/a3, reg/a4, rep/a5} {\draw[flow, dashed] (\a) -- (\b);}
    \end{tikzpicture}%
    }
    \caption{The one-week performance engagement: every phase produces an artifact.}
    \label{fig:ch9-engagement}
\end{figure}

\textbf{The expected fix.} \texttt{broadcast()} the reference side (40 GB compressed fits executor memory across the fleet---verified, not assumed), right-size shuffle partitions, and add a scheduled \texttt{ANALYZE TABLE} so AQE's statistics stop drifting. Projected: 09:15 $\rightarrow$ 05:20 with the same cluster.

\begin{pitfall}
\textbf{Promising the number before the dossier.} ``We can get you under 6 a.m.'' is a guess until stage attribution exists. Commit to the method and the evidence standard; let the measured bottleneck set the projection.
\end{pitfall}

\subsection{Risks and Mitigations}
\begin{table}[H]
\centering
\small
\begin{tabular}{@{}p{4.6cm}p{7.6cm}@{}}
\toprule
\textbf{Risk} & \textbf{Mitigation} \\
\midrule
Tuning report changes behavior subtly & Regression suite with row-for-row reconciliation \\
Salting key leaks into downstream logic & Salt columns dropped before publish; contract documented \\
Config tuned for today's volume only & Quarterly re-baseline against volume growth \\
UDF logic forked from native behavior & Unit tests pin UDF semantics against built-ins \\
Hardware scale-up masks the real fix & Plan-level fixes land first; sizing changes are separate PRs \\
\bottomrule
\end{tabular}
\caption{Week 9 scenario risks and mitigations.}
\label{tab:ch9-risks}
\end{table}

\section{Design Decision Guide}
\index{decision guide!Week 9}%
Tuning decisions compound; isolate, measure, and record each one.

\begin{table}[H]
\centering
\small
\begin{tabular}{@{}p{3.4cm}p{4.4cm}p{4.6cm}@{}}
\toprule
\textbf{Decision} & \textbf{Choose the first when} & \textbf{Choose the second when} \\
\midrule
Salting vs.\ AQE skew handling & Extreme, known hot keys & Moderate, emergent skew \\
Built-in function vs.\ pandas UDF & The function exists & Genuinely custom row logic \\
pandas UDF vs.\ Python UDF & Always, if a UDF is needed & Never the plain Python UDF \\
More partitions vs.\ more memory & Spill with healthy GC & GC pressure with small tasks \\
Tune plan vs.\ scale hardware & Always plan first & Plan is sane and hardware-bound \\
\texttt{ANALYZE TABLE} cadence & Hot, changing tables & Static reference data \\
\bottomrule
\end{tabular}
\caption{Week 9 decision guide.}
\label{tab:ch9-decisions}
\end{table}

The engagement rule from Friday applies to every row: baseline, one variable, re-measure, reconcile. A tuning change you cannot attribute is a superstition you now maintain.

\section{Operational Guidance for Week 9}
\index{operational guidance!Week 9}%
\subsection{Performance and Reliability}
Institutionalize the regression watch: p50/p95 duration per job trended weekly, with a config-change log so every performance delta has a suspect. Schedule \texttt{ANALYZE TABLE ... COMPUTE STATISTICS} on hot tables so Catalyst and AQE plan from current reality---stale statistics are the silent join-flipper. Keep the tuning experiments reproducible: the evidence table pattern from the lab becomes the team's standard for any ``make it faster'' request.

\subsection{Security and Governance Checklist}
\begin{itemize}
    \item Custom init scripts and libraries reviewed; they run with the cluster's full data access.
    \item Shuffle partitions and memory configs changed through code review, not notebook conf cells.
    \item Tuning experiments never run against production tables with write modes enabled.
    \item Broadcast targets verified against actual sizes, not assumptions (a grown dimension is a bomb).
\end{itemize}

\subsection{Troubleshooting Playbook}
\begin{table}[H]
\centering
\small
\begin{tabular}{@{}p{3.6cm}p{4.2cm}p{4.8cm}@{}}
\toprule
\textbf{Symptom} & \textbf{Likely cause} & \textbf{First action} \\
\midrule
Sudden job regression, no code change & Statistics drift flipped the join strategy & Diff the plan against last week's; refresh stats \\
Broadcast join OOMs & ``Small'' side outgrew executor memory & Measure it; drop the hint or raise memory \\
Python UDF dominates stage time & Row-at-a-time serialization & Rewrite as built-in or pandas UDF \\
Heavy GC in executor logs & Memory pressure / oversized partitions & Check spill; repartition; right-size executors \\
\bottomrule
\end{tabular}
\caption{Week 9 troubleshooting quick reference.}
\label{tab:ch9-trouble}
\end{table}

\section{Interview Q\&A}
\subsection{Concepts and Trade-offs}
\begin{enumerate}
    \item \textbf{Q: How does Catalyst decide broadcast versus sort-merge?}\\
    A: From table statistics versus \texttt{autoBroadcastJoinThreshold}; AQE revises at runtime with measured sizes. Stale statistics are why a growing dimension silently flips a job to sort-merge.
    \item \textbf{Q: When is salting better than AQE skew handling?}\\
    A: When skew is extreme and known (one key with most rows)---salting spreads it by construction, while AQE splits only after the shuffle plan exists. Salting costs a replicated small side and a two-stage aggregation.
    \item \textbf{Q: Why are pandas UDFs faster than Python UDFs?}\\
    A: Arrow columnar batches cross the JVM/Python boundary once per batch instead of once per row, and the Python side runs vectorized pandas instead of per-row interpreter calls.
    \item \textbf{Q: Your job spills heavily. Options?}\\
    A: More shuffle partitions (smaller tasks), more executor memory, a different join strategy, or less cached competition for memory. Read which stage spills before choosing.
    \item \textbf{Q: What do you check first when a job regresses with no code change?}\\
    A: The data: volume, skew, and statistics freshness. Then the plan diff---a flipped join strategy from grown input is the usual culprit.
    \item \textbf{Q: What is the fastest legitimate win in most Spark jobs?}\\
    A: Fix the join: broadcast the provably small side and prune columns/filters into the scan. Shuffle avoidance beats micro-optimization every time.
    \item \textbf{Q: How do you keep a tuned job tuned?}\\
    A: Duration trend alerts, fresh statistics (\texttt{ANALYZE TABLE}), and the config-change log---regressions are caught as data-shape changes, not discovered as SLA breaches.
    \item \textbf{Q: When is a job ``tuned enough''?}\\
    A: When the next measured improvement costs more engineering time than the compute it saves, given the SLA margin---tuning has a stopping rule, and it is economic.
\end{enumerate}

\section{Further Reading}
The Spark performance-tuning guide is the core reference \cite{sparkSQLTuning}; \textit{Spark: The Definitive Guide} gives the practitioner's tour of joins, partitions, and UDFs \cite{chambersZaharia}. The Spark documentation covers AQE configuration \cite{sparkDocs}, and the Photon page explains Databricks' execution layer \cite{databricksPhoton}.

\section*{Key Takeaways}
\begin{itemize}
    \item Tune in isolated experiments with reconciliation; confounded changes teach nothing.
    \item The five usual suspects: join strategy, shuffle partitions, skew, UDFs, and pushdown/caching.
    \item AQE handles the average case; your job is recognizing the misfit---known skew, provable broadcasts.
    \item Stay on the DataFrame API; RDDs and plain Python UDFs exit the optimizer's warranty.
    \item A performance engagement is an evidence pipeline: dossier, attribution, isolated fixes, regression, report.
\end{itemize}

\section{Exercises}
\begin{enumerate}
    \item \textbf{(Conceptual)} Explain why sort-merge joins exist at all, given broadcast joins are faster.
    \item \textbf{(Conceptual)} What information do table statistics give Catalyst, and what goes wrong when they are stale?
    \item \textbf{(Hands-on)} Complete the Thursday lab with all four experiments and the negative control.
    \item \textbf{(Hands-on)} Find one UDF in an existing notebook and replace it with built-ins; measure the delta.
    \item \textbf{(Hands-on)} Force a spill and identify it in the UI; then fix it two different ways.
    \item \textbf{(Design)} Write the shuffle-partitions rule of thumb as a decision table for 1 GB / 100 GB / 10 TB shuffles.
    \item \textbf{(Operations)} Add a scheduled \texttt{ANALYZE TABLE} to the Friday scenario's pipeline and justify its cadence.
    \item \textbf{(Architecture)} When would you move this workload to DLT (Chapter 10) instead of hand-tuned Spark? What would you lose?
\end{enumerate}
